% Adapted from Fall 2024 Final Q7 (A Pair of Schemes)

\begin{blocksection}
\question [Adapted from CS 61A Fall 2024 Final Q7 (A Pair of Schemes)] Implement \texttt{all\_pairs}, which takes a two-argument function \texttt{f} and a list \texttt{s}. It returns \lstinline{True} if \lstinline{f(x, y)} returns \lstinline{True} for every pair of adjacent elements \texttt{x} and \texttt{y} in \texttt{s}. The function \texttt{inc} returns whether \lstinline{x + 1} equals \texttt{y}.
\begin{lstlisting}
pub fn inc(x: Int, y: Int) -> Bool {
  x + 1 == y
}

pub fn all_pairs(f: fn(a, a) -> Bool, s: List(a)) -> Bool {
  case s {
    [x, y, ..rest] -> ____________ && all_pairs(f, ______________)
    _ -> True
  }
}
assert all_pairs(inc, [3, 4, 5, 6, 7, 8]) == True
assert all_pairs(inc, [3, 4, 5, 8, 7, 8]) == False
assert all_pairs(inc, [3, 4, 6, 7]) == False
assert all_pairs(inc, [3]) == True
assert all_pairs(inc, []) == True
\end{lstlisting}
\begin{solution}
\begin{lstlisting}
pub fn all_pairs(f: fn(a, a) -> Bool, s: List(a)) -> Bool {
  case s {
    [x, y, ..rest] -> f(x, y) && all_pairs(f, [y, ..rest])
    _ -> True
  }
}
\end{lstlisting}
\end{solution}
\begin{questionmeta}
%   Adjacent pairs overlap: in \lstinline{[3, 4, 5]}, the pairs are \lstinline{(3, 4)} and \lstinline{(4, 5)}, so \texttt{4} belongs to both. After checking \texttt{x} and \texttt{y}, the recursive call must start at \texttt{y}, not at the element after it. Writing \lstinline{all_pairs(f, rest)} would skip every other pair, and would incorrectly return \lstinline{True} on \lstinline{[3, 4, 6, 7]}.
%   The pattern \lstinline{[x, y, ..rest]} only matches lists with at least two elements. Every other list (empty or one element) has no pairs, so the condition holds for all of them and the \lstinline{_} case returns \lstinline{True}.
%   \lstinline{&&} only evaluates its right side when its left side is \lstinline{True}, so \texttt{all\_pairs} stops at the first pair that fails.
\end{questionmeta}
\end{blocksection}

\begin{blocksection}
\question Implement \texttt{show\_pairs}, which takes a list \texttt{s} and returns a list of all pairs of adjacent elements in \texttt{s}, in order. Each pair is a tuple.
\begin{lstlisting}
pub fn show_pairs(s: List(a)) -> List(#(a, a)) {
  case s {
    [x, y, ..rest] -> [__________, ..show_pairs(______________)]
    _ -> []
  }
}
assert show_pairs([3, 5, 7, 9, 11, 13])
  == [#(3, 5), #(5, 7), #(7, 9), #(9, 11), #(11, 13)]
assert show_pairs([1, 2]) == [#(1, 2)]
assert show_pairs([1]) == []
\end{lstlisting}
\begin{solution}
\begin{lstlisting}
pub fn show_pairs(s: List(a)) -> List(#(a, a)) {
  case s {
    [x, y, ..rest] -> [#(x, y), ..show_pairs([y, ..rest])]
    _ -> []
  }
}
\end{lstlisting}
\end{solution}
\begin{questionmeta}
  This has the same structure as \texttt{all\_pairs}, but instead of combining each pair's result with \lstinline{&&}, it collects the pairs themselves into a list. As before, the recursive call starts at \texttt{y} because adjacent pairs overlap.
  A list with $n \geq 1$ elements has $n - 1$ adjacent pairs, so the result is always one element shorter than \texttt{s}.
\end{questionmeta}
\end{blocksection}

% \begin{blocksection}
% \question What order of growth describes the time it takes to evaluate \lstinline{unpair(s)} in terms of the length of the input list \texttt{s}?
% \begin{lstlisting}
% pub fn unpair(s: List(#(a, a))) -> List(a) {
%   case s {
%     [] -> []
%     [#(x, y)] -> [x, y]
%     [#(x, _), ..rest] -> [x, ..unpair(rest)]
%   }
% }
% assert unpair(show_pairs([3, 5, 7, 9])) == [3, 5, 7, 9]
% \end{lstlisting}
% \begin{solution}
% Linear.
% \end{solution}
% \begin{questionmeta}
% %   \texttt{unpair} undoes \texttt{show\_pairs}: it keeps the first element of every pair, plus both elements of the last pair. Each call does a constant amount of work (one pattern match and one prepend) and makes at most one recursive call on a list that is one pair shorter. For a list of $n$ pairs, that is $n$ calls with constant work each, so the total is $\Theta(n)$.
% %   Prepending with \lstinline{[x, ..unpair(rest)]} takes constant time. If we had built the result with \lstinline{list.append(unpair(rest), ...)} at the end instead, each append would walk the whole result, and the total would be quadratic.
% \end{questionmeta}
% \end{blocksection}

\begin{blocksection}
\question Gleam cannot build an expression and evaluate it later the way Scheme's \texttt{eval} can, but it can build a list of \emph{delayed calls}: functions that take no arguments and make a call only when they are called.
Implement \texttt{all\_pairs\_calls}, which takes a two-argument function \texttt{f} and a list \texttt{s}. It returns one delayed call for each pair of adjacent elements \texttt{x} and \texttt{y} in \texttt{s}, which calls \lstinline{f(x, y)} when it is called. Then implement \texttt{run\_all}, which calls each delayed call in a list and returns whether all of them returned \lstinline{True}. You may use \texttt{show\_pairs}.
\begin{lstlisting}
pub fn all_pairs_calls(f: fn(a, a) -> Bool, s: List(a)) -> List(fn() -> Bool) {
  list.map(show_pairs(s), fn(pair) { ______________________________ })
}

pub fn run_all(calls: List(fn() -> Bool)) -> Bool {
  list.all(calls, ______________________)
}
assert list.length(all_pairs_calls(inc, [3, 4, 5, 6, 7, 8])) == 5
assert run_all(all_pairs_calls(inc, [3, 4, 5, 6, 7, 8])) == True
assert run_all(all_pairs_calls(inc, [3, 4, 5, 8, 7, 8])) == False
assert run_all([]) == True
\end{lstlisting}
\begin{solution}
\begin{lstlisting}
pub fn all_pairs_calls(f: fn(a, a) -> Bool, s: List(a)) -> List(fn() -> Bool) {
  list.map(show_pairs(s), fn(pair) { fn() { f(pair.0, pair.1) } })
}

pub fn run_all(calls: List(fn() -> Bool)) -> Bool {
  list.all(calls, fn(call) { call() })
}
\end{lstlisting}
\end{solution}
\begin{questionmeta}
%   Writing \lstinline{f(pair.0, pair.1)} in the first blank would call \texttt{f} right away and put a \lstinline{Bool} in the list. Wrapping it as \lstinline|fn() { f(pair.0, pair.1) }| creates a function that will make that call later. So \texttt{all\_pairs\_calls} does none of the actual checking; it only describes what to check. The fields of a tuple are accessed by position with \lstinline{.0} and \lstinline{.1}. We could also have unpacked the tuple first with \lstinline|let #(x, y) = pair| and then returned \lstinline|fn() { f(x, y) }|.
%   \texttt{run\_all} is where the calls actually happen: \lstinline{list.all} calls the function \lstinline|fn(call) { call() }| on each delayed call. Like \lstinline{&&}, \lstinline{list.all} stops at the first \lstinline{False}, so the later calls are never made.
%   Together, \lstinline{run_all(all_pairs_calls(f, s))} gives the same answer as \lstinline{all_pairs(f, s)}. The difference is when the work happens. This is the same idea as the original Scheme question, which built the expression \lstinline{(and (inc 3 4) (inc 4 5) ...)} and passed it to \texttt{eval} later.
\end{questionmeta}
\end{blocksection}
